%% _template.tex -- copy this file to start a new recipe.
%%
%% HOW TO USE
%%   1. Copy to recipes/<your_recipe>.tex (or shared/<your_subrecipe>.tex).
%%   2. Give the recipe a unique label:  [label=rec:<slug>]  (analyses)
%%                                        [label=sub:<slug>]  (shared sub-recipes)
%%   3. Fill in every field below. Delete a field ONLY if it genuinely does not
%%      apply -- prefer "None" over silence.
%%   4. \input your file from tox_manual.tex under the correct \part.
%%
%% GUIDELINE (issue #147): document reproducible WORKFLOWS, not API/source code.
%% Keep it simple, safe, clear, concise. Link to full scripts; show only the
%% fragments a user must inspect or edit.

\begin{recipe}[label=rec:CHANGE_ME]{Recipe Title}

% One or two sentences: what biological question does this answer?
\recipepurpose{TODO: state the purpose.}

% Ownership. Use \nomaintainer unless one person owns this workflow; it prints
% the team as the maintainer. \recipemaintainer{Name} names an individual.
\nomaintainer

% The data the user must have before starting.
\begin{requiredinputs}
  \item TODO: e.g. per-sample expression matrices (TPM), one file per species.
\end{requiredinputs}

% Other recipes that must be completed first. Each \dependson links by label.
\begin{prerequisites}
  \dependson{sub:CHANGE_ME}
\end{prerequisites}

% The ordered procedure. One \step per action; reference scripts, don't paste them.
\begin{workflow}
  \step{TODO: first action.}
  \step{TODO: next action.}
\end{workflow}

% The full runnable script lives in the repo; link it here.
\examplescript{path/to/example_script.sh}{%
  TODO: what this script does and what a default run produces.}

% Only the fragments the user must inspect or edit (inputs, parameters).
\begin{codefragment}[language=bash,caption={Edit the input paths and key options}]
# TODO: input data
INPUT=data/my_matrix.tsv
# TODO: key parameters
MIN_COUNT=10
\end{codefragment}

% Parameters worth deliberate choice: what each controls, and how to choose.
\begin{parameters}
  \param{--min-count}{TODO: what it controls}{TODO: when to raise/lower it}
\end{parameters}

% What the user should see at intermediate checkpoints.
\begin{expectedresults}
  TODO: describe intermediate outputs the user can check against.
\end{expectedresults}

% Traps and how to avoid them (repeatable -- add more \begin{pitfall} blocks).
\begin{pitfall}
  TODO: a common mistake and how to avoid it.
\end{pitfall}

% How to read the final output.
\begin{interpretation}
  TODO: how to interpret the results.
\end{interpretation}

% Relevant publications. Use \reference{...} or \cite{key} with references.bib.
\begin{references}
  \reference{TODO: Author et al. (year). Title. Journal.}
\end{references}

\end{recipe}
